\uuid{sopC}
\titre{Valeurs prises et support de la densité}
\auteur{Maxime Nguyen}
\organisation{AMSCC}
\datecreate{2026-11-24}
\niveau{L2}
\module{Probabilité}
\chapitre{Probabilité continue}
\sousChapitre{Densité de probabilité}
\difficulte{2}
\theme{densité, support, fonction de répartition, espérance}

\contenu{
\begin{question}
Soit $X$ une variable aléatoire de densité
$$f(x) = \frac12\left(\mathbf{1}_{[0;1]}(x) + \mathbf{1}_{[2;3]}(x)\right),$$
où $\mathbf{1}_I$ désigne l'indicatrice de l'intervalle $I$ ($1$ sur $I$, $0$ ailleurs) : $f$ vaut $\frac12$ sur $[0;1]$ et sur $[2;3]$, et $0$ ailleurs. On note $F_X$ la fonction de répartition de $X$.
Quelles affirmations sont vraies ?
\begin{answers}
\good{$\mathbb{P}(0 \leq X \leq 3) = 1$ : presque sûrement, $X$ prend ses valeurs dans $[0;3]$.}
\good{$\mathbb{P}(1 < X < 2) = 0$.}
\good{$F_X$ est constante, égale à $\frac12$, sur $[1;2]$.}
\bad{$\mathbb{P}\left(\left|X - \mathbb{E}(X)\right| < \frac12\right) > 0$.}
\bad{$\mathbb{P}(X \leq 1{,}5) = f(1{,}5) = 0$.}
\end{answers}
\begin{explanations}
A, B : $\mathbb{P}(X \in A) = \int_A f$ est nul dès que $f$ est nulle sur $A$. Les valeurs prises par $X$ sont, presque sûrement, dans l'ensemble où $f > 0$, ici $[0;1] \cup [2;3]$.

C : $F_X(t) = \int_{-\infty}^t f$ n'augmente pas là où $f$ est nulle ; elle reste plate sur $[1;2]$, à la valeur $\int_0^1 \frac12\,\mathrm{d}x = \frac12$. De même $F_X = 0$ sur $]-\infty;0]$ et $F_X = 1$ sur $[3;+\infty[$.

D : $\mathbb{E}(X) = \int_0^1 \frac{x}{2}\,\mathrm{d}x + \int_2^3 \frac{x}{2}\,\mathrm{d}x = \frac14 + \frac54 = \frac32$, donc l'événement $\left|X - \frac32\right| < \frac12$ est $\{1 < X < 2\}$, de probabilité nulle. $X$ ne prend jamais de valeur proche de son espérance.

E : $f(1{,}5) = 0$, mais $\mathbb{P}(X \leq 1{,}5) = F_X(1{,}5) = \frac12$ : toute la masse de $[0;1]$ est déjà à gauche.

Contre-sens visés : ne pas relier l'ensemble des valeurs de $X$ à l'ensemble
où la densité est non nulle ; croire que l'espérance est une valeur « typique »
de $X$ ; confondre $f(t)$ et $F_X(t)$.
\end{explanations}
\end{question}
}
